\secnumberlesssection{INTRODUCCIÓN}

La imagenología médica constituye uno de los pilares fundamentales del diagnóstico clínico moderno, proporcionando herramientas no invasivas para explorar el interior del cuerpo humano. En sus inicios, estas tecnologías se enfocaban exclusivamente en la captura de morfologías anatómicas estáticas. Sin embargo, el continuo desarrollo tecnológico ha impulsado una transición hacia la imagenología funcional, permitiendo no solo observar la estructura de los órganos, sino también cuantificar fenómenos fisiológicos dinámicos y complejos, como la biomecánica de los fluidos.

En este contexto, la Resonancia Magnética de Flujo 4D (\textit{4D Flow MRI}) se ha consolidado como una técnica avanzada que permite la adquisición y cuantificación de los campos tridimensionales de velocidad del flujo sanguíneo a lo largo del ciclo cardíaco. A diferencia de las modalidades tradicionales bidimensionales, esta tecnología captura la hemodinámica en toda su complejidad volumétrica y temporal.

No obstante, la riqueza y masividad de los datos generados por \textit{4D Flow MRI} introducen un desafío analítico y perceptivo significativo. Históricamente, la interpretación de estos campos vectoriales densos ha recaído de manera exclusiva sobre técnicas de visualización gráfica, tales como la representación por líneas de corriente (\textit{streamlines}) o mapas de calor. Este paradigma puramente visual desaprovecha por completo las capacidades innatas del sistema auditivo humano. El canal auditivo posee una sensibilidad evolutiva superior para procesar información temporal, reconocer frecuencias concurrentes y discriminar sutiles variaciones caóticas, cualidades que pueden ser sumamente útiles en el diagnóstico clínico.

Frente a esta limitación, la presente memoria propone la exploración de una dimensión analítica complementaria: la sonificación de datos. Mediante la integración de la mecánica de fluidos, el procesamiento de imágenes médicas y el procesamiento digital de señales, este trabajo plantea el diseño, implementación y validación de una metodología computacional capaz de transformar la hemodinámica adquirida por \textit{4D Flow MRI} en representaciones auditivas. La premisa central se basa en explotar la alta sensibilidad del sistema auditivo ante variaciones espectrales, estructurando una emulación del ultrasonido Doppler sobre los datos volumétricos.

Para materializar esta propuesta, el desarrollo se articuló en torno a la creación de herramientas computacionales específicas, destacando un motor acústico base y una interfaz interactiva de evaluación. Esto permitió aislar topológicamente el flujo a lo largo de su \textit{centerline} anatómico y sintetizar características acústicas coherentes. La validación empírica del modelo se llevó a cabo sobre fantomas bajo condiciones variables de morfología y caudal, incorporando finalmente técnicas de aprendizaje no supervisado para certificar el valor analítico y estadístico de los descriptores sonoros generados.

Para exponer detalladamente este proceso, el presente documento se estructura de la siguiente manera: el \textbf{Capítulo 1} detalla la formulación del problema, estableciendo los objetivos y alcances de la investigación. El \textbf{Capítulo 2} presenta el marco teórico, abordando los fundamentos físicos y computacionales necesarios para la comprensión del trabajo. El Capítulo \textbf{Capítulo 3} describe la metodología propuesta, detallando el diseño algorítmico y computacional. El Capítulo \textbf{Capítulo 4} expone los resultados obtenidos a partir de las pruebas experimentales. Finalmente, el Capítulo \textbf{Capítulo 5} sintetiza las conclusiones y presenta líneas de trabajo futuro.